\documentclass[12pt]{article}
\usepackage[utf8]{inputenc}
\usepackage[T1]{fontenc}
\usepackage{hyperref}
\title{Puyuma Language Resources for K-12 Curriculum Development}
\author{Compiled by Claude Sonnet 5.5 (claude-sonnet-5-5) for Andr\'as Kornai}
\date{}
\begin{document}
\maketitle

\section{Introduction}
Puyuma (ISO 639-3 pyu; Glottocode puyu1239) is an Austronesian language of southeastern Taiwan, classified as vulnerable by UNESCO, with most speakers being older adults. Only a few resources turned up in this search.

\section{Primary Educational Resources}
Search results mention teaching and learning materials published by Taiwan's Council of Indigenous Peoples, linked from the Wikipedia article below. I did not see the materials themselves, so I give no direct link.

\section{Linguistic and Reference Materials}
\begin{itemize}
\item Nanwang Puyuma-English Dictionary (Josiane Cauquelin, Academia Sinica, 2015, 562 pages), based on the Nanwang dialect: \url{https://www.govbooks.com.tw/books/108522}
\item Wikipedia entry on Puyuma, with links to teaching materials and an online dictionary: \url{https://www.Wikipedia.org/wiki/ISO_639:pyu}
\end{itemize}

\section{Religious and Community Resources}
None were found in this search.

\section{Digital and Audio Resources}
The Wikipedia entry above points to an online Puyuma search page on the aboriginal language dictionary site of the Indigenous Languages Research and Development Foundation. I did not see that page itself.

\section{Curriculum Development Context}
The language is endangered, and the search found no school curriculum built on it. A curriculum would have to rely on the Council of Indigenous Peoples materials and on community speakers, using Cauquelin's dictionary as a reference.
\end{document}